% ==============================================================================
% SECCIÓN 2: IDENTIFICACIÓN Y ANÁLISIS DE STAKEHOLDERS (PARTE 1 DE LA GUÍA)
% ==============================================================================

\section{Parte 1: Identificación y Caracterización de Stakeholders}

El proyecto \textbf{SIGC-CUSCO} adopta una arquitectura multientidad que articula las necesidades normativas de las entidades públicas y la agilidad requerida por el sector privado regional.

\subsection{Actores Clave del Ecosistema}
\begin{enumerate}
    \item \textbf{Municipalidades del Cusco (Sector Público):} Gerencias de Administración, Logística, Recursos Humanos y Turismo. Gestionan capacitaciones obligatorias en normativas estatales (OSCE, SIGA, SIAF, gestión ambiental) y requieren reportes fiscalizables.
    \item \textbf{Empresas del Sector Turismo (Sector Privado):} Hoteles, restaurantes y agencias de viaje. Requieren capacitar a su personal en atención al cliente, idiomas y protocolos de calidad, con acreditación digital inmediata para certificaciones y sellos de calidad.
    \item \textbf{Participantes y Ciudadanos:} Servidores públicos, trabajadores del sector turismo y estudiantes que buscan inscripción ágil y certificados infalsificables con validación QR.
    \item \textbf{Capacitadores y Expositores:} Instructores que requieren registrar asistencia de forma ágil (código QR / digital) y calificar evaluaciones sin burocracia manual.
    \item \textbf{Administradores de TI / Soporte:} Responsables de la infraestructura cloud, disponibilidad y seguridad criptográfica de las credenciales y actas.
\end{enumerate}

\subsection{Matriz de Caracterización de Stakeholders}

\begin{table}[H]
\centering
\small
\renewcommand{\arraystretch}{1.2}
\begin{tabularx}{\textwidth}{P{2.8cm}P{3.0cm}P{4.2cm}Y}
\toprule
\textbf{Stakeholder} & \textbf{Rol en el Sistema} & \textbf{Dolores Actuales (\textit{Pains})} & \textbf{Expectativas en el SIGC (\textit{Gains})} \\
\midrule
\textbf{Municipalidades (Público)} & Entidad patrocinadora y fiscalizadora. & Listas en papel, firmas adulteradas y demoras de semanas en emitir constancias. & Dashboard centralizado, control de vacantes y reportes exportables para auditoría. \\
\midrule
\textbf{Empresas Turísticas (Privado)} & Gestión de talento y competencias. & Alto costo de coordinación, pérdida de diplomas y falta de trazabilidad. & Portal institucional con marca propia, registro rápido y certificados para auditorías. \\
\midrule
\textbf{Participantes} & Beneficiario directo de la formación. & Inscripción confusa, cupos inciertos y diplomas físicos sin validación en línea. & Inscripción en 2 clics, notificaciones y certificado PDF con QR verificable. \\
\midrule
\textbf{Capacitadores} & Facilitador y evaluador académico. & Pérdida de tiempo tomando lista en papel y cálculo manual de notas. & Marcado digital de asistencia (QR), registro de notas y cierre automático de actas. \\
\midrule
\textbf{Administradores TI} & Soporte e infraestructura. & Colapso por registros masivos y riesgo de falsificación documental. & Arquitectura en la nube escalable, procesamiento asíncrono y hash SHA-256. \\
\bottomrule
\end{tabularx}
\caption{Matriz de caracterización de Stakeholders para el proyecto SIGC-CUSCO.}
\label{tab:stakeholders}
\end{table}

\subsection{Priorización de Interesados (Matriz Poder vs. Interés)}
\begin{itemize}
    \item \textbf{Gestionar de Cerca (Alto Poder / Alto Interés):} Gerentes municipales y directores de empresas turísticas (decisores de adopción).
    \item \textbf{Mantener Informados (Bajo Poder / Alto Interés):} Participantes, ciudadanos y capacitadores (usuarios operativos diarios).
    \item \textbf{Mantener Satisfechos (Alto Poder / Bajo Interés):} Entes reguladores y auditores (OSCE, Contraloría, DIRCETUR).
    \item \textbf{Monitorear (Bajo Poder / Bajo Interés):} Proveedores externos de infraestructura y público en general.
\end{itemize}
